\section{A fixed-coupling history for the rank comparison}
\label{sec:restarted-rank-history}

The rank tests compare the interacting fields with a free field over a
short interval of scales. They require an auxiliary integration history.
The coupling is held fixed on that interval, and the local additive
remainder is controlled in value. We give the density construction and
its bounds here. The calibrated history constructed earlier supplies the
endpoint and the kernels; it is still the history defining the fields.

\subsection{The additive remainder and its exact density factor}

On a finite torus $\Lambda$, write a representation's density as
\[
 R(U)=\exp\{v|\Lambda|-A(U)-E(U)\}\,G(U),
 \qquad E(U)=\sum_\alpha E_\alpha(U).
\]
Here $A$ is its path action, $G$ its forced gas with the prescribed
supports, inventories and eligibility tests, and $E_\alpha$ are real,
measurable, gauge-invariant local functions: their values agree when the local links and the relevant final flag data agree. Each $E_\alpha$ vanishes
unless its support $X_\alpha$ is healthy at the representation's threshold.
Give $X_\alpha$ an anchor $a_\alpha\in X_\alpha$. Suppose $b_\alpha\ge0$
majorizes $|E_\alpha|$ and, for the representation's support weight
$M(X)\ge1$,
\begin{equation}\label{eq:restart-anchored}
 \sum_{\alpha:a_\alpha=x}M(X_\alpha)b_\alpha\le q
 \quad(x\in\Lambda).
\end{equation}
The sums in this inequality are assumed summable. Since the anchor set
is finite, summing it gives
\[
 \sum_\alpha |E_\alpha(U)|\le |\Lambda|q.
\]
Thus the sum defining $E$ converges absolutely, is measurable and gauge
invariant, and $e^{-E}$ has the finite-volume bound $e^{|\Lambda|q}$.
Let $R^0$ be obtained by deleting just the additive terms. Its paths,
coupling, scalar, forced activities and their tests are unchanged. Then
\begin{equation}\label{eq:restart-white}
 R=e^{-E}R^0.
\end{equation}
For any finite family of source variables, the same identity holds for
its polynomial-valued density. Consequently the integrability of that
density follows from the corresponding envelope for $R^0$. The estimate
$e^{|\Lambda|q}$ is used only for finite-volume integrability; it is not
used as a uniform bound in the volume limit.

\subsection{A local multiplier keeps the expansion on the actual density}

Here is the local operation behind the construction. For a masked
additive term $E_\alpha$, insert the letter
\[
 h_\alpha=e^{-E_\alpha}-1.
\]
The identity $\prod_\alpha(1+h_\alpha)=e^{-E}$ is exact. Its infinite
version follows from absolute summability, since
$|e^{-t}-1|\le e^{|t|}|t|$ and $\sum_\alpha b_\alpha<\infty$.
If the weighted footprint estimate for these letters has allowance
$\omega$, attach them to the same component alphabet as the fluctuation
integral. The component estimate of
Section~\ref{sec:mixed-activity-prices} and the reassembly identity of
Section~\ref{sec:gaussian-reassembly} apply with this enlarged alphabet.
The old forced labels and their eligibility conditions are retained.
In particular the resulting identity integrates $e^{-E}R^0=R$.
This construction cannot be replaced by integrating $R^0$ alone.

For later use, denote by $\delta$ the weighted allowance of the
unmarked coarse letters, including the multiplier letters. The
regular-gas estimate of Section~\ref{sec:forest-average-log}, at the
purchased support exponent, supplies a masked logarithmic family
$\ell_\gamma$ with anchored cost at most $16\delta$, and
\begin{equation}\label{eq:restart-log}
 G_{\rm reg}(U)=\exp\{-\sum_\gamma\ell_\gamma(U)\}.
\end{equation}
The elementary numerical conditions used at this step include
\[
 16\delta\le1,\qquad
 e^{2\delta}-1\le1-e^{-16\delta},\qquad
 16\delta\,\mathcal R(B)\le1.
\]
Here $\mathcal R(B)$ is the connected-support counting constant in the
regular logarithm bound, and $B$ is the coarse support radius. They
are imposed on the same letter prices, not on a different auxiliary
normalization. If the integrand has holes around marked sites, the
forced-exclusion identity gives the corresponding localized logarithms
with those holes. The scalar identity without holes and these identities
with holes therefore have the same normalization.

Choose the new coarse scalar from the positive fluctuation prefactor.
If the prefactor is $C>0$, its contribution to the scalar per coarse
site is $|\Lambda_c|^{-1}\log C$. Positivity follows from the positive
coupling, the positive chart ratio, and the positive determinant of the
flat precision. In particular this logarithm does not choose a branch.
Let $A_c$ be the retained coarse path action and $\beta B_c$ the
reference action used in the fluctuation formula. Its local correction
is $\beta B_c-A_c$. Add its masked local pieces to
$\ell_\gamma$. If the former have anchored cost $\epsilon$, the
combined family has cost
\begin{equation}\label{eq:restart-newcost}
                         \epsilon+16\delta.
\end{equation}
Indeed the two nonnegative majorant sums add at every anchor. The
identity for its exponential is exactly
\[
 e^{v_c|\Lambda_c|-\beta B_c}\,G_{\rm reg}
       =e^{v_c|\Lambda_c|-A_c-E_c}.
\]
The regular/forced decomposition of the same expansion then gives
\begin{equation}\label{eq:restart-push}
                  R_c=\mathcal K_\beta R_f
\end{equation}
on every coarse configuration, including configurations with flags.
To check this last assertion, first keep the source polynomial in the
finite expansion. Its coefficient of the empty inventory equals the
unmarked scalar density. The reassembly and exclusion identities hold
coefficientwise and preserve exact flag coverage. Setting the source
family empty therefore identifies their scalar coefficients on every
configuration. There is no division by a density or by a forced
activity in this argument.

\subsection{Symmetrizing the additive family}

The newly obtained sum $E_c$ is invariant under signed lattice motions,
but an individual choice of anchored pieces need not have a symmetric
indexing. Let $H$ be the finite group of these motions on the torus.
Replace a piece $f_\alpha$ by the family
\[
 |H|^{-1}f_\alpha(gU),\qquad (g,\alpha)\in H\times\mathcal I,
\]
with the correspondingly transported support and anchor. For fixed $g$,
the anchored bound is the old bound at the inverse image of the anchor.
Multiplication by $|H|^{-1}$ and summation over $g$ preserve that bound.
Absolute summability permits interchange of the finite group sum and
the piece sum. Since the total was invariant, the new total is
$|H|^{-1}\sum_gE_c(gU)=E_c(U)$. Left multiplication in $H$ permutes the
new indices, giving the required covariance of individual addresses.
Reindex this countable family and, if necessary, add zero terms.
The gas and all forced addresses are unchanged. This produces the
locality, measurability and motion properties together with
\eqref{eq:restart-push}.

\subsection{The cost over a short interval}

Fix a block ratio $L>1$ and a width parameter $b\ge0$. For large $s$ put
\[
 y=\log\log s,\quad A_s=e^{y^3},\quad
 q_n=s^{-1/200}A_s^{4n},\quad
 \omega_n=q_nA_s,\quad \delta_n=q_nA_s^3.
\]
The two parts of the comparison interval have lengths $h_1,h_2$.
Their definitions are the ceilings
\[
 h_1=\left\lceil\frac{(b+2)y^2}{\log L}\right\rceil,
 \qquad
 h_2=\left\lceil\frac{\log(M_s^{10})}{\log L}\right\rceil,
 \qquad h=h_1+h_2,
\]
where $M_s=\lceil(\log s)^2\rceil$; the support wall is $M_s^2$.
The support-radius estimates give $\log(M_s^{10})\le y^2$ eventually.
Thus $h\le C_{L,b}y^2+2\le y^4$ eventually.
For every fixed integer $r\ge0$, uniformly for $0\le n\le h$,
\begin{equation}\label{eq:restart-decay}
 q_nA_s^r\le s^{-1/400}
\end{equation}
for all sufficiently large $s$. In fact its logarithm is at most
$-(\log s)/200+(4y^4+r)y^3$. For $y\ge1$ the positive term is at most
$(4+r)y^7$. Since $y^7e^{-y}\to0$ and $\log s=e^y$, it is at most
$(\log s)/400$ eventually, proving the claim.

The footprint price of an additive input with cost $q_n$ is at most
\[
 C_{L,a,t}(42B+1)^4q_n,
\]
where the support reserve $t$ is strictly positive, $a$ is the fixed
support exponent, and $B\le M_s^2$. All factors except the displayed
radius and $q_n$ are fixed before $s$ is increased. The wall estimates
imply $C_{L,a,t}(42M_s^2+1)^4\le A_s$ eventually. Hence the multiplier
letters cost at most $\omega_n$. The remaining unmarked letters have
their purchased local bounds; their combined logarithm is allowed cost
$16\delta_n$. The path correction costs at most $s^{-1/200}$. These
bounds close the recurrence because, as soon as $A_s\ge17$,
\begin{equation}\label{eq:restart-recurrence}
 s^{-1/200}+16\delta_n\le q_{n+1}.
\end{equation}
Indeed $s^{-1/200}\le q_n\le\delta_n$ and
$q_{n+1}=A_s\delta_n$.

For clarity, the small exponential conditions also hold uniformly
through this interval. If $0\le t\le1$, then $e^t-1\le2t$.
Consequently $e^{2\delta_n}-1\le4\delta_n$. If
$0\le\delta_n\le1/16$, the inequalities
$e^{16\delta_n}\ge1+16\delta_n$ and
$(1+16\delta_n)^{-1}\le1-4\delta_n$ give
\[
 e^{2\delta_n}-1\le1-e^{-16\delta_n}.
\]
The second rational inequality follows after multiplying by its positive
denominator: $12\delta_n-64\delta_n^2\ge0$.
The factors $\mathcal R(B)$, the insertion inflation and any fixed
comparison polynomial have at most the purchased powers of
$e^{y^2}$; multiplying \eqref{eq:restart-decay} by such factors still
gives zero in the limit. The forced-gas allowance is smaller than every
negative power of $s$, uniformly for the local scale parameter in
$[s/2,s]$. It therefore pays the support-counting factors in the
normalization-with-holes estimates. All these are eventual conditions
after $L,b$ and the finite menu have been fixed.

\subsection{Joining to the endpoint}

Choose a comparison depth $k$ and cutoff $N\ge k$. Above $k$, retain
the original history. At $k$, use its actual representation in the
chosen local presentation. Below $k$, set the coupling identically
$\beta_{N,p,k}$ and apply the preceding construction down to the lower
edge of the comparison interval. The local scale parameter lies in
$[s/2,s]$, with $s$ the clock coupling at $k$, and the endpoint bound
$|\beta_{N,p,k}-s|\le s/8$ places this constant coupling in the
purchased local range throughout the interval.

Inductively, the input at a step has cost $q_n$, the output has cost
at most $q_{n+1}$ by \eqref{eq:restart-recurrence}, and their unnormalized
densities satisfy \eqref{eq:restart-push} at that very same coupling.
Only the new entry is installed at each backward step. Thus all entries
already constructed, including the endpoint and its marked labels, are
retained. Composition proves the density identities for the entire
auxiliary history. The comparison with one or two marked sites uses the
same component alphabet and the same hole normalization. Its leading
and remainder estimates are separate from this unmarked construction;
they must be verified on these actual densities before the rank tests
can be applied.
